\begin{center}
\textsc{Gauss} an \textsc{Olbers}.\quad Göttingen, 1.~März 1827.
\end{center}

\par\noindent\hspace*{\parindent}\mbox{.\ .\ .\ .\ .\ .}\quad Meine Abhandlung, oder vielleicht richtiger, meine erste Abhandlung über die krummen Flächen habe ich vollendet; ich werde sie aber der Societät noch nicht übergeben, da doch auf die Ostermesse kein Band
\par\smallskip\noindent\hspace*{0.03\linewidth}{\footnotesize\textbf{IX.}}\hfill{\footnotesize48}\hspace*{0.14\linewidth}\par
\clearpage
\noindent\makebox[0pt][l]{\fontsize{13}{16}\selectfont 378}\hfill{\small\textbf{BRIEFWECHSEL MIT OLBERS.}}\hfill\null\par\medskip
\noindent herauskommt.
Die beiden von mir 1825 und 1826 übergebenen Abhandlungen über die biquadratischen Reste und Suppl.\ theor.\ combin.\ observ.\ sind noch nicht zu drucken angefangen.
Jene Abhandlung enthält zur \textls[120]{unmittelbaren} Benutzung in meinem künftigen Werk über die Messung eigentlich nur ein paar Sätze, nemlich 1)~was zur Berechnung des Excesses der Summe der 3 Winkel über \(180^{0}\) in einem Dreiecke auf einer \textls[120]{nicht} sphärischen Fläche, wo die Seiten kürzeste Linien sind, erforderlich ist, 2)~wie in diesem Fall der Excess auf die drei Winkel \textls[120]{ungleich} vertheilt werden muss, damit die Sinus den Seiten gegenüber proportional werden.
In praktischer Rücksicht ist dies zwar ganz unwichtig, weil in der That bei den grössten Dreiecken, die sich auf der Erde messen lassen, diese Ungleichheit in der Vertheilung unmerklich wird; aber die Würde der Wissenschaft erfordert doch, dass man die Natur dieser Ungleichheit klar begreife.
Und so kann man allerdings hier, wie öfters, ausrufen: Tantae molis erat! um dahin zu gelangen. ---
Wichtiger aber als die Auflösung dieser 2 Aufgaben ist es, dass die Abhandlung mehrere allgemeine Principien begründet, aus denen künftig, in einer speciellern Untersuchung, die Auflösung von einer Menge wichtiger Aufgaben abgeleitet werden kann.

\par\noindent\hspace*{\parindent}Ich komme noch einmal auf den im Anfange dieses Briefes erwähnten Gegenstand zurück.
Ich habe bei allen meinen Hülfstafeln und Rechnungen \textsc{Walbecks} Abplattung \(\frac{1}{302{,}78}\) zum Grunde gelegt; aber ich glaube, dass die sämmtlichen bisherigen Gradmessungen, wenn man ihre Data \textls[120]{vollständig} aufnähme, zeigen würden, dass die \textsc{Sabine}sche Abplattung \(\scriptstyle\frac{1}{2\mkern5mu 8\mkern5mu 9}\) sich beinahe ebenso gut damit würde vereinigen lassen, und es scheint mir, dass alle bisher vorhandenen Gradmessungen noch viel zu kleine Ausdehnung haben, um die Abplattung in engere Grenzen einzuschliessen.
Was gewiss ist, ist, dass die Erde ein unregelmässiger Körper ist; die Polhöhen der Örter (ganz abstrahirt von Beobachtungsfehlern) schwanken immer mehrere Secunden (vielleicht hie und da ziemlich viele Secunden) um die Werthe, die man unter Voraussetzung irgend eines regelmässigen Ellipsoids \textls[120]{berechnet}, und ebenso schwanken die wirklichen Pendellängen um die berechneten (das \textls[120]{mittlere} Schwanken der Pendellängen vielleicht \(\scriptstyle\frac{1}{3\mkern7mu 0}\) englische Linie).
Aber eben deshalb können Gradmessungen von kleiner Ausdehnung wenig zur Kenntniss der mittlern Gestalt der Erde beitragen, namentlich halte ich den lappländischen
\par\clearpage
\noindent\null\hfill{\small\textbf{ZUR HANNOVERSCHEN TRIANGULATION.}}\hfill\makebox[0pt][r]{\fontsize{13}{16}\selectfont \textbf{379}}\par\medskip
\noindent Bogen für viel zu klein.
Um so wichtiger, däucht mir, ist es, dass nach und nach alle scharfen Triangulirungen in Europa in Einen Zusammenhang gebracht werden.
Ich habe jetzt Hoffnung, die bayerischen Dreiecke mitgetheilt zu erhalten.
Wenn erst alle europäischen Sternwarten von Åbo bis Palermo und von Nicolajef bis Dublin durch Dreiecke zusammenhängen, so dass ihre relativen Lagen gegen einander mit aller Genauigkeit, die die feinsten geodätischen Messungen verschaffen können, bestimmt sind, so wird man, d.~i.\ so werden unsere Nachkommen, alles mit viel mehr Sicherheit beurtheilen \mbox{können.\ .\ .\ .\ .\ .\ .}
